% Generated draft: requires mathematical and source review.
\documentclass[11pt]{article}
\usepackage[margin=25mm]{geometry}
\usepackage{fontspec}
\setmainfont{Latin Modern Roman}
\usepackage{amsmath,amssymb,mathtools,unicode-math}
\setmathfont{Latin Modern Math}
\usepackage{enumitem,needspace,xurl,hyperref}
\setlength{\parindent}{0pt}
\setlength{\parskip}{5pt}
\setlength{\emergencystretch}{3em}
\newcommand{\N}{\mathbb{N}}
\newcommand{\Z}{\mathbb{Z}}
\newcommand{\Q}{\mathbb{Q}}
\newcommand{\R}{\mathbb{R}}
\newcommand{\C}{\mathbb{C}}
\newcommand{\F}{\mathbb{F}}
\newcommand{\A}{\mathbb{A}}
\newcommand{\PP}{\mathbb{P}}
\newcommand{\E}{\mathbb{E}}
\newcommand{\eps}{\varepsilon}
\newcommand{\dd}{\,\mathrm{d}}
\begin{document}
\section*{UnsolvedMath ID: 801042}
\textbf{Source problem code:} catalogue entry
\subsection{Definitions and mathematical statement}
The supplied question concerns hyperfiniteness of graph connectivity relations, but does not specify the ambient measurable setting or all growth quantifiers. The following is an explicit candidate Borel formulation, requiring confirmation of those missing hypotheses.

Let \(X\) be a standard Borel space and let \(G\subseteq X\times X\) be a symmetric, irreflexive Borel relation defining a locally finite graph. For \(x\in X\) and an integer \(r\geq 0\), let \(B_G(x,r)\) be the set of vertices reachable from \(x\) by a path of length at most \(r\). Define the connectivity equivalence relation by
\[
x\mathrel{E_G}y\quad\Longleftrightarrow\quad\text{there is a finite }G\text{-path from }x\text{ to }y.
\]
A Borel equivalence relation \(E\) on \(X\) is hyperfinite if there exist Borel equivalence relations \(F_n\) on \(X\), indexed by integers \(n\geq 0\), such that every \(F_n\)-class is finite,
\[
F_n\subseteq F_{n+1}\quad\text{for every }n\geq 0,
\qquad E=\bigcup_{n\geq 0}F_n.
\]
\textbf{Candidate mathematical question.} Is \(E_G\) hyperfinite for every such graph \(G\) satisfying
\[
\exists c\in(0,1)\ \exists R\in\Z_{\geq 1}\ \forall x\in X\ \forall r\in\Z_{\geq R},
\qquad |B_G(x,r)|<\exp(r^c)?
\]
Here the exponent \(c\) and radius threshold \(R\) are uniform over all vertices. This makes explicit one possible meaning of the source's phrase ``growth less than \(\exp(r^c)\) for some \(0<c<1\)''; the source does not explicitly specify these quantifiers or the Borel hypotheses.

The source calls its target ``uniform subexponential growth.'' The usual uniform subexponential ball-growth condition is instead
\[
\forall\varepsilon>0\ \exists R\in\Z_{\geq 1}\ \forall x\in X\ \forall r\in\Z_{\geq R},
\qquad |B_G(x,r)|<\exp(\varepsilon r).
\]
The displayed fixed-exponent condition is stronger than this usual condition. Consequently, whether the intended question covers all uniformly subexponential growth or only the stated fixed-exponent bounds remains unresolved.

As background, the supplied source reports hyperfiniteness for graphs of polynomial growth, attributed to Bernshteyn--Yu, and for growth less than \(\exp(r^c)\) for some sufficiently small \(0<c<1\), attributed to Grebík--Marks--Rozhoň--Shinko. It supplies neither the precise hypotheses of those results nor the meaning or numerical range of ``sufficiently small.'' These are retained as source-reported claims.
\subsection{Short English statement}
The question asks whether the connectivity relation of a graph must be hyperfinite under a uniform growth bound below \(\exp(r^c)\) for some \(0<c<1\). A candidate Borel formulation is given explicitly. The source leaves the graph setting and growth quantifiers unstated, and its use of ``uniform subexponential growth'' may indicate a broader intended question.
\subsection{Sources}
\begin{itemize}
\item Supplied local UnsolvedMath registry record, permanent ID \(801042\), exact source code \texttt{catalogue entry}, titled \emph{Open Problems in Mathematical Logic --- Problem 54}. The registry lists Barmpalias, Gao et al. as proposers, with proposed year 2026, and labels the problem open; that status is supplied metadata, not independently verified.
\item Barmpalias, Gao et al., \emph{Open Problems in Mathematical Logic} (2026), Problem 54. Supplied source URL: \url{https://arxiv.org/abs/2608.26628}. The record gives access date 2026-10-16 and extraction method \texttt{arxiv\_latex}; these are supplied provenance details, not a claim of independent access or verification.
\item Background attributions supplied in the statement: Bernshteyn--Yu for polynomial growth, and Grebík--Marks--Rozhoň--Shinko for sufficiently small exponents in bounds of the form \(\exp(r^c)\). No separate titles, bibliographic details, or links for these results are supplied.
\end{itemize}
\end{document}
